\section{Limitations and Opportunities}

\subsection{Cost Analysis}
\emph{Should you generate synthetic data, or just train longer on real data?}

The applications of \ours~lies in both paradigms---(i) low-resourced data settings such as a language model for Finnish language~\citep{luukkonen2023fingpt}, and (ii) data-rich settings such as training on the common crawl.
In the former, there is no alternative option of naively gathering more data, and hence, synthetic data is a natural solution that should outperform training on in-domain data alone. However, there is a significant interest in training language models on English, or more broadly, general web data. Is using synthetic data a viable option even in this paradigm?

Before, we dive into the feasibility of pre-training on synthetic data, we should acknowledge the results of Table~\ref{tab:general_understanding}. The TinyLlama model trained for 3 Trillion tokens also underperforms a model jointly trained on real and synthetic data. In fact, it performs quite comparably to the models that were trained for 300B tokens on just real data as well. This suggests that the ceiling for improvement by training for longer may not be that high (for a model of size 350M/1.3B parameters; larger models may benefit from training for longer).

To analyze this cost trade-off, we compare the cost of generating synthetic data, versus that of training a language model on extra data.
For our synthetic data generation experiments, we use the vLLM~\citep{kwon2023efficient} library for fast generation. In particular, we are able to generate 3M tokens per hour on a single A100 when using the Mistral-7B. Generating 85B tokens (as in our work) accounts for about 25K GPU hours. 

In comparison, on 64 A100s, we achieve a throughput of 0.5M tokens per second. Assuming training for 300B tokens, would mean 256 GPU days, accounting for about 6k GPU hours to train a single model. On the contrary, training a 13B model would take about 30K GPU hours. At the scale of training a 13B model, reducing the training cost by 3-10x can incorporate the cost overhead of training with synthetic data in a single run. 



While the cost of generating high quality data is still relatively high, two important sources of improvement impact this cost analysis.  First, if we use the Qwen-1.8B model \cite{bai2023qwen} for rephrasing, we are able to get a 3x higher token throughput. As seen in our preliminary results in Fig~\ref{fig:model_rephraser}, the model pre-trained on rephrases generated by Qwen model performs comparably to that by the Mistral model. This reduces the cost of generation by 3x. More recent work in speculative decoding~\citep{liu2023online} and optimized inference~\citep{xia2024unlocking} suggest that we can leverage another 3-5x improvement in the generation cost. Hence, indeed, even at the scale of just 1.3B parameter model training, we can already improve upon the cost of pre-training using just real data.

Two additional important advantages of synthetic data generation that could not be accounted for in the discussion above:
\begin{enumerate}
    \item The cost of synthetic data generation is a one-time investment, and we may train many models of varying scales once the data is generated.
    \item Data generation is 100\% parallelizable, whereas training requires the availability of a big cluster with fast inter-node connectivity. This is much more expensive. On the other hand, generation can be thought of as a side process that can fill in the empty GPUs in any large-scale compute cluster, and runs on single GPU machines.
\end{enumerate}

\subsection{Diversity of Synthetic Generations}
Another limitation is enforcing the diversity in the generated data. This diversity comes from both the ``style'' and the ``knowledge'' contained in the generated data. Recent works~\citep{li2023self,li2023textbooks} used a selection of topics, or scenarios to seed the model to generate novel texts.  Still, a recent study by~\citet{padmakumar-etal-2023-investigating} showed that using language models for AI-assisted writing tends to reduce content diversity, particularly when using instruction-tuned models. While we used the paradigm of rephrasing specifically to mitigate the issues pertaining to the diversity of novel content generation, it remains for future work to assess the presence (or lack of) and impact of content diversity in paraphrase models.


\section{Conclusion}
Strong language models are being pre-trained on combinations of real and synthetic data. 
Using synthetic data enables baking in desirable attributes such as fairness, bias, and style (like instruction following) directly into the data, eliminating the need to adjust the training algorithm specifically. This offers an alternative approach to aligning language models to human values. 
The recent uptick in interest around synthetic data, especially for instruction-tuning language models, is noteworthy, with concurrent researchers also leveraging it for pre-training. 
As we transition into this paradigm, understanding the properties of the data fed to our models is paramount. 
This paper aims to be a comprehensive guide on employing different synthetic style data in LLM pre-training. We delve into its significance from two vantage points: (1) In scenarios with scarce high-quality data, synthetic rephrases offer more value than mere repetition of existing data; (2) Synthetic data can be a boon for generalization on different text domains, and for generating text in styles that are underrepresented in the pre-training dataset. 
As practitioners generate synthetic data for training models, they will be faced with important and expensive design choices---(i) How important is the quality of the synthetic data generator?; (ii) How to balance real and synthetic data? (iii) When does training on synthetic data reach a point of diminishing returns in terms of epochs? This work takes a first step towards answering these questions.

Conversely, it's essential to note the inherent limitations, and opportunities with synthetic data. We highlight two limitations: (1) cost of generation is still large and requires strong LMs, and (2) enforcing the diversity in the generated data is challenging.  In this work, we leverage the natural diversity of the web to generate synthetic ``re-phrases''. This limits the model from learning new ``knowledge'' and enhances the learning process only via the provision of high-quality inputs. Whereas past work required a more intricate understanding of the blind spots of the model, potentially biasing the knowledge contained in the pre-training data distribution. Nonetheless, we demonstrate the potential of synthetic data to improve LLM training efficiency both in compute and data size. 




